\section{USAMO 2024}
        \begin{enumerate}
        	\item (\textit{USAMO 2024}) Find all integers $n \geq 3$ such that the following property holds: if we list the divisors of $n !$ in increasing order as $1=d_1<d_2<\cdots<d_k=n!$, then we have


 
 \[d_2-d_1 \leq d_3-d_2 \leq \cdots \leq d_k-d_{k-1}.\]


 

	\item (\textit{USAMO 2024}) Let $S_1, S_2, \ldots, S_{100}$ be finite sets of integers whose intersection is not empty. For each non-empty $T \subseteq\left\{S_1, S_2, \ldots, S_{100}\right\}$, the size of the intersection of the sets in $T$ is a multiple of the number of sets in $T$. What is the least possible number of elements that are in at least 50 sets?


 

	\item (\textit{USAMO 2024}) Let $m$ be a positive integer. A triangulation of a polygon is $m$-balanced if its triangles can be colored with $m$ colors in such a way that the sum of the areas of all triangles of the same color is the same for each of the $m$ colors. Find all positive integers $n$ for which there exists an $m$-balanced triangulation of a regular $n$-gon.


 Note: A triangulation of a convex polygon $\mathcal{P}$ with $n \geq 3$ sides is any partitioning of $\mathcal{P}$ into $n-2$ triangles by $n-3$ diagonals of $\mathcal{P}$ that do not intersect in the polygon's interior.


 

	\item (\textit{USAMO 2024}) Let $m$ and $n$ be positive integers. A circular necklace contains $m n$ beads, each either red or blue. It turned out that no matter how the necklace was cut into $m$ blocks of $n$ consecutive beads, each block had a distinct number of red beads. Determine, with proof, all possible values of the ordered pair $(m, n)$.


 

	\item (\textit{USAMO 2024}) Point $D$ is selected inside acute triangle $A B C$ so that $\angle D A C=$ $\angle A C B$ and $\angle B D C=90^{\circ}+\angle B A C$. Point $E$ is chosen on ray $B D$ so that $A E=E C$. Let $M$ be the midpoint of $B C$.
Show that line $A B$ is tangent to the circumcircle of triangle $B E M$.


 

	\item (\textit{USAMO 2024}) Let $n>2$ be an integer and let $\ell \in\{1,2, \ldots, n\}$. A collection $A_1, \ldots, A_k$ of (not necessarily distinct) subsets of $\{1,2, \ldots, n\}$ is called $\ell$-large if $\left|A_i\right| \geq \ell$ for all $1 \leq i \leq k$. Find, in terms of $n$ and $\ell$, the largest real number $c$ such that the inequality

 \[\sum_{i=1}^k \sum_{j=1}^k x_i x_j \frac{\left|A_i \cap A_j\right|^2}{\left|A_i\right| \cdot\left|A_j\right|} \geq c\left(\sum_{i=1}^k x_i\right)^2\]
holds for all positive integers $k$, all nonnegative real numbers $x_1, \ldots, x_k$, and all $\ell$-large collections $A_1, \ldots, A_k$ of subsets of $\{1,2, \ldots, n\}$.


 Note: For a finite set $S,|S|$ denotes the number of elements in $S$.


 

\end{enumerate}